% === preamble.tex ===
% Все пакеты и настройки для сборника задач

% --- Кодировка и язык ---
\usepackage[T2A]{fontenc}
\usepackage[utf8]{inputenc}
\usepackage[russian]{babel}

% --- Математика ---
\usepackage{amsmath, amssymb, amsthm}
\usepackage{mathtools}

% --- Оформление задач (красивые рамки) ---
\usepackage[most]{tcolorbox}

% --- Ссылки и оглавление ---
\usepackage{hyperref}
\hypersetup{
    colorlinks=true,
    linkcolor=blue,
    urlcolor=cyan,
    pdfauthor={IronSamuray},
    pdftitle={Сборник задач по физике 10-11 класс}
}
\pdfinfo{
    /Author (IronSamuray)
    /Title (Сборник задач по физике 10-11 класс)
    /Subject (Все права защищены)
    /Creator (IronSamuray)
}
% --- Макет страницы ---
\usepackage{geometry}
\geometry{a4paper, top=2cm, bottom=2cm, left=2cm, right=2cm}

% --- Свои команды для удобства ===

% Рейтинг задачи (как на Codeforces: 500-3000)
\newcommand{\rating}[1]{%
    \textcolor{gray}{\small\textbf{[★ #1]}}%
}

% Окружение для задачи
\newtcolorbox{problem}[2][]{%
    colback=white,
    colframe=black!75,
    fonttitle=\bfseries,
    title=Задача #2,
    #1 % дополнительные параметры
}

% Окружение для решения (можно скрывать)
\newif\ifsolutions
\solutionstrue % Закомментировать, чтобы скрыть решения

\newenvironment{solution}{%
    \ifsolutions%
    \begin{tcolorbox}[colback=yellow!5, colframe=orange!50, title=💡 Решение]%
}{%
    \ifsolutions%
    \end{tcolorbox}%
    \fi%
}

% Теги для задач (ЕГЭ, олимпиада, база и т.д.)
\newcommand{\tags}[1]{%
    {\footnotesize\textcolor{gray}{\textit{#1}}}%
}